\section{Discussion}\label{sec:discussion}

\subsection{What the benchmark tests}
Problem S verifies the anisotropic, active, nearly incompressible solid
against nine independent codes. Problems F and C add the fluid while keeping
a solid-only reference: the difference between each FSI solution and its AR25
counterpart measures the effect of the blood, and is itself a quantity other
codes can reproduce. Problem F exercises the added-mass instability under
passive loading ($\rho_{\mathrm f}/\rho_{\mathrm s}\approx1$) and the
formation of the filling vortex; Problem C exercises active, fibre-driven
deformation that pushes the fluid out, with large wall thickening and twist
that challenge the mesh motion.

\subsection{Pressure-driven vs.\ flow-driven ports}\label{sec:disc:dilemma}
\TODO{Discuss F-p vs.\ F-Q: in F-Q the cavity volume is imposed, the pressure
acts as a Lagrange multiplier, and Dirichlet--Neumann partitioned coupling
fails without a volume constraint or a Robin interface condition
\citep{Kuttler2006}. Show how the Robin condition in solids4foam handles it
and the cost in coupling iterations. Monolithic codes are unaffected.}

\subsection{Recommended numerical settings}
\TODO{From the robustness map: which stabilisation factors give mesh-converging
results without artificial stiffening; which coupling acceleration and tolerance
give the best accuracy per cost; minimum mesh and $\dt$ for each QoI. Which QoIs
converge quickly (global: $V$, $\bar p$, displacements) and which slowly
(local: $\Delta p_{\mathrm{IVPD}}$, vortex metrics, $E_k$).}

\subsection{Limitations}
Filling and contraction are posed separately, not as a closed cycle;
idealised geometry without papillary muscles or trabeculae; no valve
leaflets; Newtonian, laminar blood; activation prescribed rather than
electrophysiology-driven. The AR25 ventricle is large ($V_0=\SI{168}{mL}$
unloaded), more representative of a dilated than a healthy LV, and the AR25
pressure waveform reaches systolic levels. These are deliberate choices that
keep the benchmark unambiguous and comparable with AR25; a closed-cycle
benchmark with lumped-parameter circulation is a natural next step.

\subsection{Recommendations for reporting}
\Cref{app:checklist} lists the minimum information that, in our view, a
ventricular FSI study should report to be reproducible.
